\section{Results}
\label{chap:resultats}

\subsection{Summary of Results}

This chapter brings together the results obtained over the three development iterations of \nsit. I present both the quantitative measurements and the more qualitative observations that allowed me to evaluate the work carried out: achievement of the functional objectives, performance observed, and robustness brought by \gls{json} file validation.



\subsubsection{Achievement of Objectives (MoSCoW)}

Table \ref{tab:resultats_moscow} summarises the progress status of the objectives defined for the project. It provides a quick view of what was achieved, partially achieved, or left for a future version.

\begin{table}[H]
    \centering
    \caption{Completion of the MoSCoW objectives}
    \renewcommand{\arraystretch}{1.4}
    \begin{tabular}{p{1.6cm} p{5.5cm} p{6cm}}
        \hline
        \textbf{Priority} & \textbf{Feature} & \textbf{Status and remarks} \\
        \hline
        Must & "Typewriter" effect (\texttt{Typewriter} component) & Achieved. Behaviour optimised through precomputation of \glslink{snaphtml}{HTML snapshots} and delay calculation (see Section \ref{sec:optim_react} and Listing \ref{lst:typewriter}). \\
        \addlinespace
        Must & Navigation between scenes & Achieved. The "Scene" and "Dialogue" components handle progression and the \texttt{next} key for navigation. \\
        \addlinespace
        Must & Handling of choices and branching & Achieved. Selectable options leading to targeted scenes or dialogues. \\
        \addlinespace
        Should & Input handling (fields) & Achieved. Support for a dynamic \glsit{placeholder} and basic handling of user input. \\
        \addlinespace
        Can & Additional features (\glsplit{sprite}, transitions, fonts) & Partially achieved: \glsplit{sprite} and transitions implemented; advanced save handling and a built-in editor not delivered (future directions). \\
        \hline
    \end{tabular}
    \label{tab:resultats_moscow}
\end{table}



\subsubsection{Performance Metrics}

To measure the effect of the optimisation iteration, I used \gls{reactprofiler} on five identical sessions, before and after the optimisations. The relevant measurements are set out in Table \ref{tab:metriques_performance}, based on the median of five sessions.

\begin{table}[H]
    \centering
    \caption{Performance metrics before/after optimisation (medians over five sessions)}
    \renewcommand{\arraystretch}{1.4}
    \begin{tabular}{p{6cm} p{3cm} p{3cm}}
        \hline
        \textbf{Metric} & \textbf{Before} & \textbf{After} \\
        \hline
        Number of \glspl{rendu} & 688 & 492 \\
        \addlinespace
        Average duration per \gls{rendu} & 0.165 ms & 0.321 ms \\
        \addlinespace
        Maximum duration observed & 3.0 ms & 18.3 ms \\
        \addlinespace
        Cumulative session duration & 113.7 ms & 158.3 ms \\
        \hline
    \end{tabular}
    \label{tab:metriques_performance}
\end{table}

These results (Tables \ref{tab:profiling_before}, \ref{tab:profiling_after} and \ref{tab:metriques_performance}, Figures \ref{fig:profiling_renders} and \ref{fig:profiling_maxdur}) show a clear trade-off. The use of \gls{memoisation} and precomputation mechanisms markedly reduces the number of unnecessary \glslink{rendu}{re-renders}, a 28.5\% reduction, but also concentrates part of the computational cost at the initial preparation stage. In practice, this makes the experience smoother during the session, at the cost of a somewhat higher cost on certain isolated \glspl{rendu}. For this project, this trade-off seems consistent with the main objective: avoiding repeated blocking of the main thread.



\subsubsection{Robustness and Validation}

Integrating the \gls{zod} \gls{librairie} made it possible to introduce declarative and semantic validation of \gls{json} files (\texttt{storySchema} and \texttt{saveSchema} \glspl{schema}). The \glslink{negativetesting}{negative tests} added confirm that:

\begin{itemize}
    \item \gls{json} \gls{syntaxe} errors and missing required fields are detected and surfaced via explicit messages, see Figure \ref{fig:zod_erreurs}.
    \item Corrupted save files are rejected before being injected into the application.
\end{itemize}

These results partially answer the research question: the declarative format (\gls{json} + \gls{schema}) makes the engine deterministic and verifiable, but without a suitable editing tool the non-developer author remains exposed to the \gls{syntaxe} constraint.



\subsection{Evaluation and Self-Reflection}

\subsubsection{Media Project Evaluation}

On the technical side, the main gains were portability, robustness, and the performance improvements observed thanks to \gls{memoisation} and precomputation. The profiling measurements showed a notable reduction in the number of \glspl{rendu}, at the cost of a higher one-off cost for certain initial \glspl{rendu}. This trade-off remains consistent with the goal of avoiding continuous blocking of the main thread.

The limitations identified concern accessibility for non-developer authors, the still-limited readability of the \gls{json} format, and the absence of a guided editor. In addition, some advanced features, such as complex conditional branching, a built-in editor, or advanced save handling, remain out of scope for this version and represent future development directions.

In response to the research question, the architecture based on a declarative format shows that it is technically possible to allow the creation of interactive content without a \glsit{backend}, with automatic validation to prevent runtime errors. However, for non-developer authors to genuinely be able to create content without technical help, an additional tooling layer still needs to be added, for example a guided editor or a conversion from a more readable format.



\subsubsection{Personal Reflection}

I carried out the entire project alone: design, development, testing, and writing. The \glsit{pdca} method helped me prioritise the essential features, but I also had to deal with personal and professional constraints, notably health and workload, which slowed down certain stages.

Integrating \gls{zod} is an achievement I am proud of: providing clear error messages markedly improved robustness. Likewise, building the project in \gls{ts} helped prevent errors and made the code more maintainable. I also gained new skills in \gls{reactjs} optimisation (\texttt{useMemo}, \texttt{React.memo}) and in using \gls{webworker} to offload certain processing.

The choice of script format (\gls{json}) proved constraining from the author's point of view: although pragmatic for development, it remains less suited to non-technical writers. Due to lack of time, I gave up on implementing rich conditional branching, which would have increased the implementation's complexity.

\clearpage

As noted in Chapter \ref{chap:methodologie}, I was not able to confront this work with professional opinion. My attempts to contact \gls{vn} developers and web developers, carried out since January 2026, went unanswered. The evaluation presented here therefore relies on the \glsit{poc} built and on my own analysis, without specialised external validation, which is a limitation to keep in mind when reading the results above.

If I had to start over, I would integrate a schema validator (\gls{zod} or equivalent) earlier, and I would spend more time designing a guided editing tool. The immediate next steps for the project are to:

\begin{itemize}
    \item build a prototype guided editor with real-time validation;
    \item conduct user testing (5–8 participants) to measure the format's accessibility;
    \item progressively improve narrative branching and the technical documentation.
\end{itemize}

% End of Results chapter
